% =============================================================================
% SECTION 2: SYSTEM ARCHITECTURE OVERVIEW
% =============================================================================

\section{System Architecture Overview}
\label{sec:arch}

Figure~\ref{fig:arch} summarizes the proposed architecture.
The design follows three architectural patterns common in modern AI systems~\cite{lewis2020rag}:

\textbf{(1) Retrieval-Augmented Generation (RAG)}---agents retrieve relevant facts from the KG before generating outputs, grounding decisions in structured knowledge rather than parametric memory alone.

\textbf{(2) Tool-Augmented Agents}---LLM agents invoke deterministic tools (optimizer, exporter) for tasks requiring provable correctness~\cite{schick2023toolformer}.

\textbf{(3) Blackboard Architecture}---the KG serves as a shared workspace where specialized agents read inputs, write inferences, and coordinate asynchronously.

The core design principle is that the system is \emph{KG-centered}: extracted plan facts and inferred construction intent are continuously grounded against a structured graph of rules, components, and historical builds.

\begin{figure}[ht]
\centering
\begin{tikzpicture}[
  node distance=6mm and 8mm,
  box/.style={draw, rounded corners, align=center, inner sep=4pt, font=\scriptsize, minimum width=18mm},
  arrow/.style={-Latex, thick}
]
\node[box] (ingest) {Plan Ingest\\(PDF/CAD/Images)};
\node[box, right=of ingest] (extract) {Geometry + Text\\Extraction};
\node[box, right=of extract] (facts) {Plan Facts\\(walls, dims)};
\node[box, below=8mm of facts] (kg) {\textbf{Knowledge Graph}\\Rules \& History};
\node[box, right=of facts] (agents) {Agent\\Orchestrator};
\node[box, right=of agents] (outputs) {Outputs\\BOM, Cuts,\\Instructions};

\draw[arrow] (ingest) -- (extract);
\draw[arrow] (extract) -- (facts);
\draw[arrow] (facts) -- (agents);
\draw[arrow] (kg) -- (agents);
\draw[arrow] (facts) -- (kg);
\draw[arrow] (agents) -- (outputs);
\end{tikzpicture}
\caption{KG-centered architecture: plan facts are extracted, normalized, stored with provenance, and used by agents to generate compliant, optimized outputs.}
\label{fig:arch}
\end{figure}

\subsection{Key Modules}

\textbf{(1) Plan Ingest + Extraction.}
Convert drawings into machine-usable primitives:
vector geometry (lines, polylines), detected room/wall boundaries, and text tokens (dimensions, callouts, schedules)~\cite{pizarro2022floorplan_survey}.
The current implementation supports PDF (via PyMuPDF~\cite{pymupdf2024}), DXF (via ezdxf~\cite{ezdxf2024}), and DWG (via LibreDWG conversion).
For raster images, YOLOv8 object detection~\cite{jocher2023yolov8} and EasyOCR~\cite{easyocr2024} provide robust extraction.

\textbf{(2) Plan Facts Model.}
Normalize extraction into a consistent intermediate representation:
\emph{WallSegments}, \emph{Openings}, \emph{Headers}, \emph{Notes}, \emph{Schedules}, and \emph{Constraints}
with explicit provenance back to the sheet, region, and confidence score.

\textbf{(3) Knowledge Graph.}
Store rules, materials, codes, and historical knowledge (Section~\ref{sec:kg}).

\textbf{(4) Agent Orchestrator.}
A multi-agent workflow that queries the KG, performs reasoning and validation,
and calls deterministic optimization routines for cutting and procurement.

\textbf{(5) Export + Integration.}
Produce both human-usable artifacts (PDF/HTML summaries, build steps)
and machine-readable exports (CSV/JSON/EDI-style mappings) suitable for lumber yard systems.

\subsection{Data Flow}

The system operates in five sequential phases:

\textbf{Phase 1: Extraction.}
Plan documents are parsed into vector geometry, raster images, and text.
YOLOv8 detects symbols; PyMuPDF extracts text and vectors from PDFs; ezdxf handles CAD formats.
Output: raw geometry primitives and text tokens with bounding boxes.

\textbf{Phase 2: Normalization.}
Raw extraction is normalized into \texttt{PlanFact} nodes in the KG, each with source provenance (sheet, region, confidence).
Cross-validation checks dimension consistency and flags ambiguities.

\textbf{Phase 3: Inference.}
Agents query the KG for plan facts and rules, then write inferred assemblies (\texttt{AssemblyIntent}, \texttt{Component}) back to the KG.
Each inference links to triggering facts and applied rules.

\textbf{Phase 4: Optimization.}
The Procurement Agent selects stock items and invokes the cut-stock optimizer (ILP solver) for waste minimization.
Results are written as \texttt{CutPiece} nodes linked to source components.

\textbf{Phase 5: Export.}
Final outputs (BOM, cut lists, instructions) are generated by querying the KG and formatting results.
Every output line item carries provenance links back through phases 4→3→2→1.
